\documentclass{report}
\usepackage{amsfonts, amsmath}
\newcommand\R{{\mathbb R}}
\newcommand\Q{{\mathbb Q}}
\newcommand\N{{\mathbb N}}
\newcommand\ve{\varepsilon}
\newcommand\Co{{\mathcal C}}
\pagestyle{empty}
\begin{document}
\large
\centerline{\Huge Fisica Matematica II}
\renewcommand{\thefootnote}{ } 
\footnote{\sl Avete 2:30 ore di tempo. Ogni esercizio
vale dieci punti. La sufficienza si ottiene con un punteggio $\geq$
18. Solo le risposte chiaramente giustificate saranno prese in
considerazione.}
\renewcommand{\thefootnote}{\arabic{footnote}} 
\vskip.1cm
\centerline{\Large Appello 04-09-2003}
\vskip1cm
\begin{enumerate}
\item Si determini una soluzione dell'equazione
\[
\begin{split}
&yu_y+x^2\ln y\; u_x=1\\
&u(s,e)=s\quad s\in\R
\end{split}
\]
e se ne discuta dominio e unicit\`a.
\item Si consideri l'equazione
\[
\Delta u(x)=f(x)\quad x\in\R^3,
\]
con $sup_{x\in\R^3} (|x|^4+1)|f(x)|<\infty$. Si mostri che $u\in L^\infty(\R^3)$.
\item Si consideri l'equazione
\[
\begin{split}
&u_{tt}-u_{xx}=\sin\lambda t \quad x\in[0,1]\\
&u(x,0)=u_t(x,0)=0\\
&u(0,t)=u(1,t)=0.
\end{split}
\]
Per quali $\lambda$ esiste una costante $a_\lambda>0$ tale che
$|u(x,t)|\leq a_\lambda$ per ogni $(x,t)\in[0,1]\times\R$ ?
\item Si determini la soluzione del problema
\[
\begin{split}
&u_{xx}+xe^{-t}=u_t\quad x\in[0,1],\;t>0\\
&u(x,0)=x(1-x)\quad x\in[0,1]\\
&u(t,0)=u(t,1)=0\quad t>0.
\end{split}
\]
\end{enumerate}
\end{document}
